\begin{figure}[htbp]
\centering
\begingroup
\setlength{\fboxsep}{6pt}
\renewcommand{\arraystretch}{1.15}

\begin{minipage}{0.95\textwidth}
\centering

\fcolorbox{guindaipn}{white}{%
    \begin{minipage}{0.87\linewidth}
    \centering
    \textbf{Entradas principales}\\[-2pt]
    {\small
    Datos históricos de la WNBA, registros \textit{play-by-play}, alineaciones, posesiones,
    estadísticas individuales y disponibilidad de jugadoras.}
    \end{minipage}
}

\\[4pt]
{\color{guindaipn}\Large$\Downarrow$}
\\[4pt]

\fcolorbox{guindaipn}{white}{%
    \begin{minipage}{0.87\linewidth}
    \centering
    \textbf{Preprocesamiento y validación de la información}\\[-2pt]
    {\small
    Limpieza, integración de fuentes,  identificadores, reconstrucción
    de quintetos y verificación de consistencia temporal.}
    \end{minipage}
}

\\[4pt]
{\color{guindaipn}\Large$\Downarrow$}
\\[4pt]

\fcolorbox{guindaipn}{white}{%
    \begin{minipage}{0.87\linewidth}
    \centering
    \textbf{Fase 1: caracterización de jugadoras}\\[-2pt]
    {\small
    Selección de variables, análisis de estadísticas históricas, construcción de perfiles
    funcionales y comparación entre jugadoras rivales.}
    \end{minipage}
}

\\[4pt]
{\color{guindaipn}\Large$\Downarrow$}
\\[4pt]

\fcolorbox{guindaipn}{white}{%
    \begin{minipage}{0.87\linewidth}
    \centering
    \textbf{Fase 2: evaluación de quintetos candidatos}\\[-2pt]
    {\small
    Identificación del quinteto rival, generación de alineaciones factibles del equipo propio
    y estimación preliminar del rendimiento esperado de cada combinación.}
    \end{minipage}
}

\\[4pt]
{\color{guindaipn}\Large$\Downarrow$}
\\[4pt]

\fcolorbox{guindaipn}{white}{%
    \begin{minipage}{0.87\linewidth}
    \centering
    \textbf{Optimización de la alineación}\\[-2pt]
    {\small
    Comparación de quintetos factibles mediante una función objetivo sujeta a restricciones,
    como número de jugadoras, disponibilidad y condiciones de sustitución.}
    \end{minipage}
}

\\[4pt]
{\color{guindaipn}\Large$\Downarrow$}
\\[4pt]

\fcolorbox{guindaipn}{white}{%
    \begin{minipage}{0.87\linewidth}
    \centering
    \textbf{Salidas del modelo}\\[-2pt]
    {\small
    Quinteto recomendado, ordenamiento de alternativas, estimación de rendimiento esperado
    y medidas de evaluación para el análisis posterior.}
    \end{minipage}
}

\end{minipage}
\endgroup
\caption{Resumen del procedimiento preliminar propuesto para el desarrollo del modelo de optimización de alineaciones.}
\label{fig:procedimiento_preliminar_modelo}
\end{figure}